\section{Calibration}
\label{sec:calibration}

\subsection{Parameter Values}
\label{sec:calibration_params}

The model is calibrated at quarterly frequency following
\citet{Gali2008}, Chapter~3. Table~\ref{tab:calibration} lists the
deep parameters together with their sources.

\begin{table}[t]
  \centering
  \caption{Calibrated parameters (quarterly).}
  \label{tab:calibration}
  \begin{tabular}{lcll}
    \toprule
    Parameter & Value & Description & Source \\
    \midrule
    $\beta$      & $0.99$   & discount factor                 & \citet{Gali2008}, Table~3.1 \\
    $\sigma$     & $1.00$   & inverse intertemporal elasticity & \citet{Gali2008}, Ch.~3 \\
    $\varphi$    & $1.00$   & Frisch labour disutility curvature & \citet{Gali2008}, Ch.~3 \\
    $\theta$     & $0.75$   & Calvo price stickiness          & \citet{Gali2008}, Ch.~3 \\
    $\phi_{\pi}$ & $1.50$   & Taylor rule inflation reaction  & \citet{Taylor1993} \\
    $\phi_{y}$   & $0.125$  & Taylor rule output reaction     & \citet{Gali2008}, Ch.~3 \\
    $\rho_{a}$   & $0.90$   & TFP shock persistence           & \citet{Gali2008}, Ch.~3 \\
    $\rho_{b}$   & $0.50$   & preference shock persistence    & \citet{Gali2008}, Ch.~3 \\
    $\rho_{m}$   & $0.50$   & monetary shock persistence      & \citet{Gali2008}, Ch.~3 \\
    $\sigma_{m}$ & $0.25\%$ & monetary innovation std.        & \citet{Gali2008}, Ch.~3 \\
    \bottomrule
  \end{tabular}
\end{table}

The discount factor and shock persistences follow the canonical
calibration in \citet{Gali2008}. The intertemporal elasticity is set to
unity, the textbook log-utility benchmark. The Frisch labour-supply
elasticity is also set to unity, an intermediate value consistent with
the micro estimates surveyed by \citet{ChettyEtAl2011}. The Calvo
parameter $\theta=0.75$ implies that the expected duration of a fixed
price is four quarters, matching the median duration documented by
\citet{NakamuraSteinsson2008} for the U.S.\ Consumer Price Index. The
Taylor rule coefficients follow \citet{Taylor1993} for inflation and
\citet{Gali2008} for output.

\subsection{Implied Composite Coefficient}
\label{sec:calibration_kappa}

The deep parameters imply a Phillips-curve slope of
\begin{equation}
  \kappa
  =\frac{(1-\theta)(1-\beta\theta)}{\theta}\,(\sigma+\varphi)
  =\frac{(0.25)(0.2575)}{0.75}\,(2.00)
  \approx 0.1717,
  \label{eq:kappa}
\end{equation}
which is the value used throughout the rest of the paper. Steady-state
analysis is trivial in log-deviation form: every variable equals zero by
construction, and substituting the analytical steady state into the
equilibrium conditions yields residuals identically equal to zero.
